\section{The exact four singular losses at the infinite-period boundary}
The canonical four marked ES--Fable columns remain those of
FC15a--32 \cite{ES488}. We calculate their original four-dimensional
weighted conductor, with cutoff $D=3$, both original source orders,
all Gamma factors and centres as in WCF2--6 \cite{WCF,DLMF18Source}.
The earlier programme proof PZ1--13, PZX1--6 and PZU1--11
\cite{PeriodZeroOriginal} already establishes the two exceptional
unit phases, the exact second inverse-period order of each zero-value
factor, the bound $v\le32$, and the finite domain with $v=0$.
Those are inputs, not new results here. The calculation below evaluates
the entire inverse matrix and all four exterior losses on the first
of those actual phase loci. It does not choose the original unit phase.

\subsection{Exact two-variable coefficients of the original factors}
Write $z=u^{-1}$ and use $\xi$ only for the original exponential
variable. Retain
\[
 \beta=2(\gamma^2-\delta^2),\quad
 \eta=(\delta^2+\gamma^2)^2,\quad
 B=\beta^2-9\eta=-5\gamma^4-26\delta^2\gamma^2-5\delta^4<0,
 \quad h=9/B.
 \tag{PIN1}
\]
On the actual phase locus $a=\bar a\ne0$, the complete matrix
$U=aI$ cancels exactly between $U^{-1}$ and $U$; its value is not
assigned. For $q=\zeta^{-j}$, $1\le j\le4$, set
\[
 T_j(z)=\mathcal R(z)^{-1}
       \operatorname{diag}(q,q^2,q^3,q^4)\mathcal R(z),\qquad
 p(\xi)=V^{-1}(e^{w_a\xi})_{a=1}^4,\quad w_a=\rho_a-\tfrac12.
 \tag{PIN2}
\]
Thus $f_j(\xi,z)=e^\xi Q_0(VT_j(z)p(\xi))$ is precisely OFM1,
including its original matrix order. Evaluation at the four roots of
$w^4+\beta w^2+\eta$ proves, for every coefficient column $c$,
\[
 Q_0(Vc)=4i\delta\gamma\,c^{\mathsf T}Qc,\quad
 Q=\begin{pmatrix}0&0&1&0\\0&-1&0&\beta\\
 1&0&-\beta&0\\0&\beta&0&\eta-\beta^2\end{pmatrix}.
 \tag{PIN3}
\]
Indeed the two products are $p(x)p(-x)-p(y)p(-y)$,
$x=\delta+i\gamma$, $y=\delta-i\gamma$;
$x^2-y^2=4i\delta\gamma$,
$x^4-y^4=-\beta(x^2-y^2)$ and
$x^6-y^6=(\beta^2-\eta)(x^2-y^2)$ give every displayed entry.
For $0\le k\le3$, $[\xi^k]p(\xi)=e_k/k!$, since evaluation
of the polynomial $w^k$ gives exactly that derivative.

Put $F_j=e^{-\xi}f_j/(4i\delta\gamma)$. The complete contributing
coefficients are
\[
 \begin{split}
 L_j=[z^2]F_j(0,z)&=-\frac{\beta^2B}{81}
                         (q^3+2q^2+3q-1),\\
 M_j=[z\xi]F_j(\xi,z)&=\frac{\beta^2}{9}
                         q^3(q-1)(2q^2+2q+1),\\
 [\xi^0z^0]F_j=[\xi^1z^0]F_j
   =[\xi^2z^0]F_j=[\xi^3z^0]F_j&=0.
 \end{split}\tag{PIN4}
\]
The first coefficient is PZX2 divided by $4i\delta\gamma$.
Here is a finite coefficient proof of the other assertions and a
reproducible formula retaining every term. Use PCL2's full $R_n$, and
\[
 \begin{aligned}
 S_0&=R_0^{-1},\quad
 S_n=-S_0\sum_{r=1}^nR_rS_{n-r},\quad
 T_{j,n}=\sum_{r=0}^n S_rD_jR_{n-r},\\
 [\xi^kz^n]F_j&=
 \sum_{a=0}^k\sum_{r=0}^n
 \frac{(T_{j,r}e_a)^{\mathsf T}Q(T_{j,n-r}e_{k-a})}
      {a!(k-a)!},\quad k\le3.
 \end{aligned}
 \tag{PIN5}
\]
Multiplication of $\mathcal R^{-1}\mathcal R=I$ proves the inverse
recurrence, and substitution of PIN3 proves the last sum. Inserting
PCL4 and PZX1's full $R_1$ gives PIN4 after reduction by
$1+q+q^2+q^3+q^4=0$. Thus this is an identity over the original
parameter field. It uses neither sampled coordinates nor assigned
moments. The supplied exact checker verifies every entry of this
finite calculation and all four roots by cyclotomic reduction.

The parity identity PCL22 gives
$F_j(-\xi,-z)=F_j(\xi,z)$. Together with PIN4 and the entire
coefficient series it gives, coefficient by coefficient in $\xi$,
\[
 [\xi^0]F_j=L_jz^2+O(z^4),\quad
 [\xi^1]F_j=M_jz+O(z^3),\quad
 [\xi^2]F_j=O(z^2),\quad [\xi^3]F_j=O(z).
 \tag{PIN6}
\]
All remainders are convergent Taylor remainders for the fixed original
quartet and unit. None replaces an original coefficient.

\subsection{The full inverse in both original Gamma norms}
Set
\[
 C=125(4\delta\gamma)^4\beta^8B^4/81^4>0,\qquad
 \mathcal F(\xi,z)=\prod_{j=1}^4(4i\delta\gamma F_j(\xi,z))
                         =e^{-4\xi}E_{A,u}(\xi).
 \tag{PIN7}
\]
Expanding all four factors in PIN6 and reducing fifth roots gives
\[
 [\xi^k]\mathcal F(\xi,z)
 =C a_kz^{8-k}+O(z^{10-k}),\quad 0\le k\le3,
 \qquad (a_0,a_1,a_2,a_3)=(1,-2h,2h^2,-h^3).
 \tag{PIN8}
\]
For an explicit certificate, the leading coefficient before the
factor $(4i\delta\gamma)^4$ is
$\sum_{|I|=k}\prod_{j\in I}M_j\prod_{j\notin I}L_j$.
Its four values are respectively
$125\beta^8B^4/81^4$,
$-250\beta^8B^3/9^7$,
$250\beta^8B^2/9^6$,
$-125\beta^8B/9^5$.
Terms using a second or third $\xi$ coefficient in a single factor
have at least two extra powers of $z$ by PIN6. This proves the errors
as well as every leading coefficient in PIN8.

PZU2--4 supplies the exact nonzero finite domain $|u|\ge R_A$,
on which $v=0$. Write $\omega(t)=-i\arctan t$ with its original
branch and
$g(t,z)=\mathcal F(\omega(t),z)$, $b(t,z)=1/g(t,z)$.
The original WCF6 matrix and its entire inverse at cutoff three are
\[
 \begin{gathered}
 A_{3,s}(z)=R_s^{-1}T_3(g)R_s,\quad
 A_{3,s}(z)^{-1}=R_s^{-1}T_3(b)R_s,\quad
 \\
 R_s=\operatorname{diag}(\rho_0,\rho_1,\rho_2,\rho_3),\quad
 \rho_m=\sqrt{m!/(s/2)_m}.
 \end{gathered}
 \tag{PIN9}
\]
The centres remain $a_{\rm amp}/2$ and $a_{\rm amp}/2-4$;
the original mass $(2\pi)^{s/2}$ cancels only in the orthonormal
matrix as proved in WCF6, and remains in both spaces.
Every finite-period entry, including its whole Taylor tail, is
defined by the exact recursion
\[
 b_0=g_0^{-1},\qquad
 b_n=-g_0^{-1}\sum_{k=1}^n g_kb_{n-k}\quad(1\le n\le3).
 \tag{PIN10}
\]
Since $\omega(t)=-it+it^3/3+O(t^5)$, PIN8 and this recursion give
\[
 b_n(z)=C^{-1}d_nz^{-8-n}+O(z^{-6-n}),\quad
 (d_0,d_1,d_2,d_3)=(1,-2ih,-2h^2,ih^3).
 \tag{PIN11}
\]
The cubic term of $\omega$ contributes two higher powers of $z$,
and is included in the exact recursion; it has not been discarded
from the finite-period map.

Let $\sigma_1\ge\sigma_2\ge\sigma_3\ge\sigma_4$ be the inverse
singular values in precisely these original norms. In the exterior
matrix of order $k$, each entry is a minor with row set $I$ and
column set $J$, and has pole at most
$8k+\sum_{j\in J}j-\sum_{i\in I}i$ by PIN11. This is maximized
uniquely at $I=\{0,\ldots,k-1\}$,
$J=\{4-k,\ldots,3\}$, with value $8k+k(4-k)$.
The four leading corner determinants of $T_3(d)$ are
\[
 d_3=ih^3,\quad d_2^2-d_1d_3=2h^4,\quad
 d_1^3-2d_1d_2+d_3=ih^3,\quad 1.
 \tag{PIN12}
\]
They are nonzero since $h\ne0$. Multiplying the exterior matrix
by $z^{8k+k(4-k)}$ therefore yields a matrix with one nonzero
limiting entry. Its Euclidean operator norm is the absolute value
of that entry. Norm continuity proves the complete four limits
\[
 \begin{split}
 \lim_{|u|\to\infty}|u|^{-11}\|A_{3,s}^{-1}\|
       &=\frac{\rho_3}{\rho_0}\frac{|h|^3}{C},\\
 \lim_{|u|\to\infty}|u|^{-20}\|\wedge^2 A_{3,s}^{-1}\|
       &=\frac{\rho_2\rho_3}{\rho_0\rho_1}\frac{2|h|^4}{C^2},\\
 \lim_{|u|\to\infty}|u|^{-27}\|\wedge^3 A_{3,s}^{-1}\|
       &=\frac{\rho_3}{\rho_0}\frac{|h|^3}{C^3},\\
 \lim_{|u|\to\infty}|u|^{-32}|\det A_{3,s}^{-1}|
       &=C^{-4}.
 \end{split}\tag{PIN13}
\]
These limits hold on every direction of the punctured finite-period
domain, for each unchanged source order $s=1,a_{\rm amp}$.
The equality $\|\wedge^k A^{-1}\|=\prod_{j=1}^k\sigma_j$
follows by applying the exterior functor to a unitary singular value
decomposition; its diagonal entries are all products of $k$ singular
values. Dividing successive nonzero limits in PIN13 proves
\[
 \begin{aligned}
 \sigma_1&\sim\frac{\rho_3|h|^3}{C}|u|^{11},&
 \sigma_2&\sim\frac{2\rho_2|h|}{\rho_1 C}|u|^9,\\
 \sigma_3&\sim\frac{\rho_1}{2\rho_2 C|h|}|u|^7,&
 \sigma_4&\sim\frac{1}{\rho_3C|h|^3}|u|^5,
 \end{aligned}\tag{PIN14}
\]
where $\rho_0=1$ has only been evaluated in this last display.

Finally $T_A\Phi_*q=\Psi_*q$ continues to hold for every one of
the four marked columns and every signed affine state at these
finite periods by FC31 and OFM17. Thus the entire inverse just
calculated acts on their actual original target, retaining $K^{-1}$
and every interpolation and Gamma coordinate map. At $z=0$, PIN6
implies that all four entries $g_0,\ldots,g_3$ vanish, so the fixed
old $D=3$, $v=0$ matrix has zero boundary value. PZU4 proves it is
invertible for every finite $|u|\ge R_A$. The boundary degeneration
and its exact growth therefore do not assert a zero in the finite
annulus $R_*\le|u|<R_A$, or an RH conclusion. That annulus still
has the actual finite exceptional set PZU5; it has not been replaced
by a freely chosen polynomial or a substituted unit phase.
